\documentclass[10pt,a4paper]{amsart}
\usepackage[margin=19mm]{geometry}
\usepackage{amsmath,amssymb,mathtools}
\usepackage[scr=rsfs,cal=euler]{mathalfa}
\usepackage{xfrac}
\usepackage[unicode,hidelinks,bookmarks=false]{hyperref}
\newcommand{\ie}{i.e.\ }
\newcommand{\cf}{cf.\ }
\newcommand{\resp}{respectively\ }
\newcommand{\Subsection}{\subsection}
\providecommand{\nobreakdash}{\nobreak\hbox{-}}
\setlength{\emergencystretch}{2em}
\multlinegap0pt
\usepackage{bm}
\usepackage{bbm}
\usepackage{dsfont}
\usepackage{xspace}
\usepackage{cancel}
\usepackage{mathtools}
\usepackage[shortlabels]{enumitem}
\usepackage{amsthm}
\makeatletter
\@ifundefined{theorem}{\newtheorem{theorem}{Theorem}}{}
\@ifundefined{lemma}{\newtheorem{lemma}[theorem]{Lemma}}{}
\makeatother
\newcommand{\E}{\mathbb{E}}
\newcommand{\R}{\mathbb{R}}
\newcommand{\ind}{\mathbbm{1}}
\title[Average partial effect estimation using double machine learn]{Average partial effect estimation using double machine learning}
\author{Harvey Klyne, Rajen D. Shah}
\date{Source: arXiv:2308.09207v3 (CC BY 4.0)}
\begin{document}
\maketitle

\noindent\emph{Source and licence.} Average partial effect estimation using double machine learning. Version arXiv:2308.09207v3 (CC BY 4.0), distributed under the Creative Commons Attribution 4.0 International licence, as recorded in the arXiv licence metadata for this exact version and shown on the abstract page https://arxiv.org/abs/2308.09207v3. Full rights notice: licenses/arxiv-2308.09207-CC-BY-4.0.txt.\par\smallskip

\noindent\textit{Editorial scope.} Counted unit: the Lemma whose source label is \texttt{lem:scoreboundaryconvergent} in the article (source0.tex lines 2395--2398, source label \texttt{lem:scoreboundaryconvergent}), delivered together with the article's own complete original proof of it (source0.tex lines 2399--2444, one \texttt{proof} environment of 46 lines). No mathematics is written, repaired or re-derived: statement and proof are the article's own text line for line. Required context retained uncounted: lem:pconvergent (lines 2303--2306); lem:boundaryconvergent (lines 2333--2341). Every definition, display and label of those lines is retained unchanged. Omitted from this package and not redistributed: 1--2302, 2307--2332, 2342--2394, 2445--2749. Nothing inside a retained excerpt is omitted or altered: every retained body is delivered exactly as the article's lines are written, so no omission receipt of any kind accompanies this package. Citations: the retained excerpts carry no citation command at all, so no bibliography entry of the article is transcribed and no reference is redistributed. Reference adaptation: none was needed and none was made; every reference inside the delivered unit resolves inside it, so no unresolved reference or citation is produced. Printed slips kept literal: no printed word, symbol, hypothesis or number of the counted statement or of its proof was altered. Layout and class: the article's own class is \texttt{imsart} with packages including inputenc, amsmath, amsfonts, amssymb, amsthm, enumitem, natbib, caption, babel, fontenc, graphicx, mathrsfs, bbm, array; the wrapper is the lane's pinned \texttt{amsart} editorial wrapper at 10pt on A4, which restates the theorem-like environments the retained text needs and adds the article's own macro definitions that the retained text needs (\texttt{\textbackslash E}, \texttt{\textbackslash R}, \texttt{\textbackslash ind}), with extra wrapper packages (bm, bbm, dsfont, xspace, cancel, mathtools, enumitem). The delivered numbering is the unit's own. Assets: no retained span contains a figure environment, a graphics-inclusion command, a PDF-page inclusion, a table or a TikZ picture; no image file is copied into this package, the package declares no asset dependency and no image file is redistributed with it. Licence: version arXiv:2308.09207v3 of this article is distributed under the Creative Commons Attribution 4.0 International licence, as recorded in the arXiv licence metadata for this exact version and shown on the abstract page https://arxiv.org/abs/2308.09207v3. A full rights notice is delivered at licenses/arxiv-2308.09207-CC-BY-4.0.txt. Characters: every retained excerpt is pure ASCII, so the delivered engine, XeLaTeX with its default Latin Modern text font, reports no missing character.
	\begin{lemma}
		\label{lem:pconvergent}
		If $p$ is a twice differentiable density function on $\R$ and $\sup_{x\in\R} |\rho'(x)| \leq C$, then $p(x) \to 0$ as $|x|\to\infty$.
	\end{lemma}

	\begin{lemma}
		\label{lem:boundaryconvergent}
		Let $f:\R\to\R$ be a continuous function. Then at least one of the following holds.
		\begin{enumerate}[label=(\alph*)]
			\item There exists a sequence $a_n \to \infty$ such that $f(a_n) \to 0$.
			\item There exists $A\in\R$ and $\epsilon > 0$ such that $f(x) >\epsilon$ for all $x\geq A$, and in particular $\int_A^\infty f(x)\;dx = \infty$.
			\item There exists $A\in\R$ and $\epsilon > 0$ such that $f(x) <-\epsilon $ for all $x\geq A$, and in particular $\int_A^\infty (-f(x))\;dx = \infty$.
		\end{enumerate}
	\end{lemma}

	\begin{lemma}
		\label{lem:scoreboundaryconvergent}
		Let $p$ be a twice differentiable density function on $\R$ with score $\rho$ defined everywhere, and let $k$ be a non-negative integer. If $\sup_{x\in\R} |\rho'(x)| \leq C$ and $\E[\rho^{2k}(X)]<\infty$, then there exist sequences $a_n \to -\infty$ and $b_n \to \infty$ such that $\rho^{2k+1}(a_n)p(a_n) \to 0$ and $\rho^{2k+1}(b_n)p(b_n) \to 0$.% $\rho^{2k+1}(b_n)p(b_n) - \rho^{2k+1}(a_n)p(a_n) \to 0$.
	\end{lemma}

	\begin{proof}
		Write $f(x)=\rho^{2k+1}(x)p(x)$. Since $f$ is continuous, we may apply Lemma~\ref{lem:boundaryconvergent} to both $f$ and $x \mapsto f(-x)$ to conclude that either the statement of the lemma holds, or one of the following hold for some $B \in \R$ and $\epsilon > 0$:
		\begin{enumerate}[label=(\alph*)]
			\item $f(x) > \epsilon$ for all $x \geq B$,
			\item $f(x) < - \epsilon$ for all $x \geq B$
		\end{enumerate}
		or one of the above with $x \geq B$ replaced with $x \leq B$. Let us suppose for a contradiction that (a) occurs (the other cases are similar), so in particular
		\begin{equation}
			\int_{B}^\infty f(x)\; dx = \infty. \label{eqn:rho2k+1contradictioninfinite}
		\end{equation}
		
		%Therefore it suffices to show that for any $B\in \R$, both $\int_{-\infty}^B f(x)\; dx$ and $\int_{B}^\infty f(x)\; dx$ are finite. Suppose, for contradiction, that there exists $\epsilon > 0$ and $B\in\R$ such that $f(x)\geq \epsilon$ for all $x\geq B$ and
		%	\begin{equation}
			%		\int_{B}^\infty f(x)\; dx = \infty; \label{eqn:rho2k+1contradictioninfinite}
			%	\end{equation}
		%	the other cases are similar.
		
		If $k=0$, then \begin{equation*}
			\int_B^\infty f(x)\;dx = \int_B^\infty p'(x)\;dx = \lim_{b \to \infty} \int_B^b p'(x)\;dx = \lim_{b\to\infty} p(b) - p(B);
		\end{equation*}
		here the penultimate equality follows from monotone convergence and the final equality follows from the fundamental theorem of calculus.
		By Lemma~\ref{lem:pconvergent} however, this is finite, a contradiction. If instead $k\geq 1$, then for any $b\geq B$ we have that \begin{align}
			\rho^{2k}(b)p(b)-\rho^{2k}(B)p(B) &= \int_B^b \rho^{2k}(x)p'(x)\; dx + 2k\int_B^b \rho'(x)\rho^{2k-1}(x)p(x)\;dx.\nonumber\\
			&= \int_B^b \rho^{2k+1}(x)p(x)\; dx + 2k\int_B^b \rho'(x)\rho^{2k-1}(x)p(x)\;dx.\label{eqn:rho2k+1contradictionfinite}
		\end{align}
		We will take the limit as $b\to\infty$. Since $\rho^{2k}(x)p(x)$ is non-negative and we have that $\E[\rho^{2k}(X)]<\infty$, we can choose an increasing sequence $b_n\to\infty$ satisfying $\rho^{2k}(b_n)p(b_n) \leq 1$ for every $n$. 
		
		Note that for each $n$ and for every $x\in\R$,
		\begin{equation*}
			\big|\ind_{[B, b_n]}(x)\rho'(x) \rho^{2k-1}(x)p(x)\big|
			\leq C |\rho(x)|^{2k-1}p(x).
		\end{equation*}
		By Jensen's inequality, $\E\big[|\rho(X)|^{2k-1}\big]< \infty$. Thus, by dominated convergence theorem,
		\begin{align*}
			\lim_{n \to \infty} 2k \int_B^{b_n} \rho'(x) \rho^{2k-1}(x) p(x)\;dx &= 2k \int_B^{\infty} \rho'(x) \rho^{2k-1}(x) p(x)\;dx\\
			&\leq 2kC \int_B^{\infty} |\rho(x)|^{2k-1} p(x)\;dx\\
			&\leq 2kC \;\E\big[|\rho(X)|^{2k-1}\big]< \infty.
		\end{align*}
		Now \eqref{eqn:rho2k+1contradictionfinite} implies that \begin{equation*}
			\lim_{n\to\infty} \int_B^{b_n}  f(x)\; dx < \infty.
		\end{equation*}
		But we assumed that $f(x) \geq \epsilon > 0$ for all $x\geq B$, so for each fixed $x\in\R$ the integrand $\ind_{[B,b_n]}(x) f(x) $ is increasing as a function of $n$. Therefore monotone convergence implies that \begin{equation*}
			\int_B^{\infty}  f(x)\; dx < \infty,
		\end{equation*}
		contradicting \eqref{eqn:rho2k+1contradictioninfinite}.
	\end{proof}

\end{document}
